\documentclass[a4paper]{article}
\usepackage{amsfonts}
\usepackage{amsmath}
\usepackage{amssymb}
\usepackage{graphicx}     
\usepackage{cancel}

% Command definities van frequent gebruikte commandos
\newcommand{\N}{\mathbb{N}}
\newcommand{\Q}{\mathbb{Q}}
\newcommand{\R}{\mathbb{R}}
\newcommand{\C}{\mathbb{C}}
\newcommand{\Bgsin}{\operatorname{Bgsin}}
\newcommand{\Bgcos}{\operatorname{Bgcos}}
\newcommand{\Bgtan}{\operatorname{Bgtan}}
\newcommand{\cis}{\operatorname{cis}}
\newcommand{\Limit}{\lim\limits_}

\begin{document}
  
\noindent \large Auteur: Wout Vanderborght\\
\noindent \large Studierichting: Informatica\\
\noindent \large Oefeningenreeks 3H nr. 1.4\\

\medskip

\normalsize

Schrijf de volgende vergelijking als poolcoördinaten, en schrijf deze waar \\

mogelijk als $r = f(\theta)$: $9x^2 + y^2 = 4y$ \\

$ \Leftrightarrow 9r^2\cos^2(\theta) + r^2\sin^2(\theta) = 4r\sin(\theta)$ \\

$ \Leftrightarrow r^2(9\cos^2(\theta) + \sin^2(\theta)) = 4r\sin(\theta)$ \\

$ \Leftrightarrow r^2(8\cos^2(\theta) + 1) = 4r\sin(\theta)$ \\

$ \Leftrightarrow r(8\cos^2(\theta) + 1) = 4\sin(\theta)$ of $r = 0$ \\

$ \Leftrightarrow r = \dfrac{4\sin(\theta)}{8\cos^2(\theta) + 1}$ of $r = 0$ \\

\begin{thebibliography}{999}
%\bibitem{a} Referentie
\end{thebibliography}
\end{document} 
